\chapter*{پیشگفتار}
\addcontentsline{toc}{chapter}{پیشگفتار}

\section*{پیدایش یک فریم‌ورک امن کوانتومی}

\begin{sloppypar}
تحول از ایده‌های نظری در توزیع کلید کوانتومی (\lr{Quantum Key Distribution, QKD}) به سامانه‌هایی که بتوان آن‌ها را در سطح پیاده‌سازی، اعتبارسنجی، بهینه‌سازی و نهایتاً استقرار شبکه‌ای بررسی کرد، مستلزم عبور از یک شکاف بنیادی میان دو جهان متفاوت است: از یک سو، جهان فیزیک کوانتومی و نظریهٔ اطلاعات که با صورت‌بندی‌های انتزاعی، کران‌های امنیتی، فروض ایده‌آل و مدل‌های تحلیلی سروکار دارد؛ و از سوی دیگر، جهان مهندسی سامانه و نرم‌افزار که در آن، جزئیات اجرایی، قیود محاسباتی، ناپایداری‌های سخت‌افزاری، ساختار داده، همزمانی، تکرارپذیری و اعتبارسنجی، نقش تعیین‌کننده در قابل‌اعتماد بودن یک نتیجه ایفا می‌کنند. در عمل، بخش عمده‌ای از دشواری QKD نه در فهم اصول اولیهٔ آن، بلکه در پر کردن دقیق، روشمند و قابل دفاع همین فاصله نهفته است.

مستندی که پیش رو قرار دارد، حاصل تلاشی برای پاسخ دادن به این نیاز است: طراحی، توسعه، سازمان‌دهی و اعتبارسنجی یک \textbf{فریم‌ورک شبیه‌سازی QKD} که بتواند به‌عنوان یک بستر پژوهشی جامع، دقیق، توسعه‌پذیر و مهندسی‌شده، رفتار سامانه‌های \lr{Discrete-Variable QKD} را از سطح تولید پالس نوری تا سطح پردازش امنیتی و حتی مسیریابی شبکه‌ای مدل‌سازی کند. این فریم‌ورک صرفاً یک ابزار محاسبه‌گر نرخ کلید نیست، بلکه تلاشی است برای بازسازی یک خط لولهٔ کامل از رخدادهای فیزیکی، آماری، الگوریتمی و معماری که در یک سامانهٔ واقعی QKD حضور دارند. در این معنا، پروژهٔ حاضر نه‌تنها به مدل‌سازی یک پروتکل، بلکه به ساخت یک زیرساخت تحقیقاتی برای مطالعهٔ خانواده‌ای از سناریوهای QKD می‌پردازد.

انگیزهٔ اصلی شکل‌گیری این فریم‌ورک از یک مشاهدهٔ نسبتاً روشن اما عمیق آغاز شد: بسیاری از ابزارها و مدل‌های موجود، اگرچه برای محاسبهٔ تقریبی نرخ کلید، بازتولید چند نمودار مرجع، یا اجرای سناریوهای آموزشی مفیدند، اما برای پژوهش‌های جدی‌تر با محدودیت‌های بنیادی مواجه می‌شوند. بخشی از این محدودیت‌ها از ساده‌سازی بیش از حد لایهٔ فیزیکی ناشی می‌شود؛ برای مثال، مدل‌سازی ایده‌آل یا ناقص آشکارسازها، نادیده‌گرفتن آبشارهای زمان مرده، صرف‌نظر کردن از \lr{afterpulsing}، یا فروکاستن نویز به یک ثابت ساده. بخش دیگر از محدودیت‌ها به لایهٔ امنیتی بازمی‌گردد؛ جایی که برخی ابزارها تنها یک اثبات امنیتی خاص را به‌صورت سخت‌کدگذاری‌شده پیاده‌سازی می‌کنند و در نتیجه، امکان مقایسهٔ دقیق و منصفانهٔ چند اثبات بر روی یک موتور فیزیکی واحد را از بین می‌برند. افزون بر این، در بسیاری از شبیه‌سازهای موجود، تحلیل‌ها در رژیم مجانبی انجام می‌شوند؛ در حالی که در سامانه‌های واقعی، کلید محدود، نوسانات آماری و بودجهٔ امنیتی، بخش جدایی‌ناپذیر مسئله هستند.

فریم‌ورک حاضر برای رفع این مصالحه‌ها توسعه یافت. هدف آن ایجاد محیطی یکپارچه بود که در آن، پژوهشگر بتواند بدون قربانی کردن دقت فیزیکی، انعطاف‌پذیری رمزنگاری، شفافیت معماری و قابلیت توسعه، سناریوهای متنوع QKD را مطالعه کند. از همین رو، طراحی این سامانه از ابتدا بر اساس چند اصل هدایت‌کننده شکل گرفت: تفکیک روشن لایه‌ها، ماژولار بودن، قابلیت اعتبارسنجی تحلیلی، پشتیبانی از سناریوهای \lr{Finite-key}، توانایی مقایسهٔ چند اثبات امنیتی، و امکان توسعه به سطح شبکه‌های چندگرهی. این اصول نه صرفاً ترجیحات مهندسی، بلکه پاسخ‌هایی ضروری به ماهیت میان‌رشته‌ای و حساس QKD بودند.

این پروژه در عین آنکه از نظر مفهومی بر مبانی استاندارد QKD تکیه دارد، به‌صورت خاص برای پاسخ به نیازهای پژوهش تجربی و شبیه‌سازی دقیق طراحی شده است. به همین دلیل، تمرکز آن تنها بر «درست بودن» فرمول‌ها نیست، بلکه بر «درست بودن رفتار سامانه در تعامل اجزا» نیز هست. در عمل، یک اثبات امنیتی دقیق نیز اگر بر داده‌هایی سوار شود که به‌دلیل مدل‌سازی ناقص آشکارساز یا مدیریت نادرست رویدادهای همزمان مخدوش شده‌اند، خروجی معناداری نخواهد داشت. به همین ترتیب، یک موتور فیزیکی دقیق نیز اگر نتواند داده‌های خود را به‌درستی در اختیار لایهٔ پساپردازش قرار دهد، ارزش پژوهشی محدودی خواهد داشت. از این رو، این فریم‌ورک بر پیوستگی دقیق و مهندسی‌شدهٔ کل زنجیره تأکید دارد.

به بیان دیگر، این فریم‌ورک تلاشی است برای آنکه شبیه‌سازی QKD از سطح یک اسکریپت محاسباتی مجزا فراتر رود و به یک محیط پژوهشی قابل اتکا تبدیل شود؛ محیطی که بتواند برای طراحی آزمایش، مقایسهٔ اثبات‌ها، تحلیل حساسیت، بهینه‌سازی پارامتر، بررسی سناریوهای چندگرهی و تولید نتایج بازتولیدپذیر به کار رود. این مستند نیز با همین نگاه تدوین شده است: نه فقط برای توضیح آنچه پیاده‌سازی شده، بلکه برای روشن‌سازی منطق علمی، معماری نرم‌افزاری و مسیر اعتبارسنجی‌ای که این فریم‌ورک بر آن بنا شده است.

\section*{فلسفه معماری}

در طراحی این فریم‌ورک، معماری نرم‌افزار تنها یک مسئلهٔ سازمان‌دهی کد نبوده، بلکه بخشی از منطق علمی پروژه به‌شمار آمده است. هنگامی که یک سامانه قرار است به‌طور همزمان فیزیک منبع، تلفات کانال، دینامیک آشکارساز، آمارداده‌های کلید محدود، اثبات‌های امنیتی متنوع، بهینه‌سازی عددی و حتی رفتار شبکه‌ای را در خود جای دهد، هرگونه آمیختگی نامنظم میان این اجزا به‌سرعت به پیچیدگی کنترل‌ناپذیر، دشواری در اعتبارسنجی و کاهش قابلیت اطمینان منجر خواهد شد. به همین دلیل، فلسفهٔ معماری این پروژه بر پایهٔ \textbf{تفکیک مسئولیت‌ها، ماژولار بودن، صراحت در قراردادهای داده، و استقلال نسبی لایه‌ها} شکل گرفته است.

در هستهٔ این طراحی، تمایز روشنی میان دو ناحیهٔ اصلی برقرار شده است: نخست، \textbf{لایهٔ فیزیکی} که مسئول مدل‌سازی آن چیزی است که در سطح منبع، کانال و آشکارساز رخ می‌دهد؛ و دوم، \textbf{لایهٔ پساپردازش و امنیت} که وظیفهٔ تفسیر داده‌های مشاهده‌شده، غربال‌گری پایه، تخمین پارامتر، اعمال اثبات‌های امنیتی و محاسبهٔ نرخ کلید امن را بر عهده دارد. این جداسازی دو مزیت مهم ایجاد می‌کند. از یک سو، امکان آن را فراهم می‌سازد که مدل‌های فیزیکی بدون دست‌کاری مستقیم در لایهٔ اثبات‌ها توسعه یابند؛ و از سوی دیگر، به پژوهشگر اجازه می‌دهد چندین اثبات امنیتی را بر روی دقیقاً یک داده‌مولد فیزیکی مشترک اجرا کند. به این ترتیب، مقایسه‌ها نه تحت تأثیر تفاوت پیاده‌سازی‌های پراکنده، بلکه بر پایهٔ یک بستر همگن و کنترل‌شده انجام می‌شوند.

یکی از عناصر شاخص این معماری، \textbf{موتور توزیع چندکارگری (\lr{Multi-Worker Dispatch Engine})} است. در شبیه‌سازی‌های رویدادمحور QKD، مقیاس مسئله به‌سرعت بزرگ می‌شود. زمانی که تحلیل در رژیم \lr{Finite-key}، روی تعداد قابل توجهی پالس، با مدل‌های واقع‌گرایانه و مجموعه‌ای از پارامترها انجام می‌شود، هزینهٔ محاسباتی اجرای ترتیبی به‌شدت افزایش می‌یابد. در چنین شرایطی، استفاده از پردازش موازی دیگر یک بهینه‌سازی لوکس نیست، بلکه یک الزام عملی است. با این حال، موازی‌سازی در این حوزه با چالش‌هایی همراه است که در بسیاری از سامانه‌های ساده نادیده گرفته می‌شوند. مهم‌ترین این چالش‌ها، حفظ \textbf{یکپارچگی هویت پالس‌ها و سازگاری رویدادها} در محیطی است که اقلام کاری روی کارگرهای مختلف، با زمان‌بندی‌های متفاوت و احتمالاً با تأخیرهای ناهمگون پردازش می‌شوند.

برای پاسخ به این مسئله، معماری فریم‌ورک از مکانیزمی مبتنی بر \lr{\texttt{work item}}های مستقل و خودبسنده بهره می‌برد. هر قلم کاری، اطلاعات کافی دربارهٔ پالس، پایه، شدت، حالت و زمینهٔ آماری خود را به همراه دارد و به‌گونه‌ای تعریف می‌شود که وابستگی پنهان یا شکننده به وضعیت اجرایی دیگر بخش‌ها نداشته باشد. این تصمیم معماری باعث می‌شود که نتایج به ترتیب زمان‌بندی کارگرها وابسته نشوند و پردازش موازی به تغییر خاموش در معنای داده‌ها منجر نگردد. در یک سامانهٔ رمزنگاری یا شبه‌رمزنگاری، این موضوع اهمیتی فراتر از کارایی دارد؛ زیرا هرگونه ناسازگاری خاموش در هویت پالس‌ها یا تعلق رویدادها می‌تواند مستقیماً به تفسیر نادرست آمار امنیتی بینجامد.

از منظر طراحی داده، این فریم‌ورک تلاش می‌کند تا حد ممکن از قراردادهای صریح، ساختارهای دادهٔ مشخص و مسیرهای پردازش قابل رهگیری استفاده کند. این امر صرفاً برای زیبایی کد یا سهولت توسعه نیست؛ بلکه برای آن است که یک نتیجهٔ عددی نهایی، در صورت نیاز، قابل ردیابی به مراحل پیشین خود باشد. پژوهشگر باید بتواند بفهمد یک نرخ کلید یا یک نرخ خطا چگونه از میان لایه‌های مختلف شکل گرفته است، چه عواملی بر آن اثر گذاشته‌اند، و در صورت بروز انحراف، منشأ آن در کدام بخش قرار دارد. در اینجا، معماری خوب عملاً نقش ابزار تبیین علمی را نیز ایفا می‌کند.

فلسفهٔ معماری پروژه همچنین بر \textbf{توسعه‌پذیری کنترل‌شده} تأکید دارد. در سامانه‌های پژوهشی، نیازها ایستا نیستند. ممکن است در نسخه‌های بعدی، مدل جدیدی از منبع نوری، یک تقریب بهتر برای \lr{afterpulsing}، یک اثبات امنیتی متفاوت، یا یک سیاست تازه برای مدیریت شبکه مورد نیاز باشد. اگر افزودن هر قابلیت جدید مستلزم بازنویسی چندین ماژول و شکستن فرضیات قبلی باشد، سامانه خیلی زود به نقطهٔ فرسودگی می‌رسد. از این رو، طراحی کلاس‌های پایه، لایه‌های انتزاعی و رابط‌های داخلی به‌گونه‌ای انجام شده است که افزودن جزء جدید، عمدتاً مستلزم پیاده‌سازی همان جزء باشد، نه بازآرایی کل سیستم.

در نهایت، فلسفهٔ این معماری را می‌توان چنین خلاصه کرد: \textit{دقت علمی تنها زمانی ارزشمند است که در یک ساختار مهندسی‌شده، شفاف، تکرارپذیر و قابل توسعه استقرار یابد.} این فریم‌ورک تلاش کرده است دقیقاً چنین بستری را فراهم کند؛ بستری که در آن، شبیه‌سازی QKD از یک محاسبهٔ موردی و محدود فراتر رود و به یک دارایی پژوهشی و نرم‌افزاری پایدار تبدیل شود.

\section*{فرآیند جامع اعتبارسنجی}

هیچ فریم‌ورک شبیه‌سازی، صرف‌نظر از پیچیدگی ظاهری، تنوع قابلیت‌ها یا ظرافت معماری، بدون یک فرایند اعتبارسنجی دقیق و چندلایه، شایستهٔ اعتماد علمی نخواهد بود. این گزاره در مورد سامانه‌های QKD اهمیتی دوچندان دارد؛ زیرا خطا در چنین محیطی تنها یک اختلاف عددی ساده نیست، بلکه می‌تواند به تفسیر نادرست امنیت، خوش‌بینی کاذب نسبت به نرخ کلید، یا مقایسه‌های گمراه‌کننده میان مدل‌ها و اثبات‌ها منجر شود. از همین رو، اعتبارسنجی در این پروژه نه یک مرحلهٔ پسینی و تکمیلی، بلکه بخشی مرکزی از فرایند توسعه بوده است.

رویکرد اعتبارسنجی این فریم‌ورک بر چند اصل بنا شده است. نخست آنکه اعتبارسنجی نباید به سطح تست واحد توابع محدود شود، هرچند این تست‌ها لازم‌اند. دوم آنکه هر لایه باید تا حد امکان در برابر یک مرجع مستقل سنجیده شود: لایهٔ فیزیک در برابر روابط تحلیلی و پیش‌بینی‌های فرم‌بسته، لایهٔ امنیتی در برابر خواص مورد انتظار اثبات‌ها و رفتار مرزی آن‌ها، و لایهٔ شبکه در برابر منطق ساختاری سناریوهای مسیریابی و تأمین کلید. سوم آنکه اعتبارسنجی باید هم در سطح محلی و هم در سطح برهم‌کنش میان لایه‌ها انجام شود؛ زیرا بسیاری از خطاهای مهم نه در خود اجزا، بلکه در مرزهای میان آن‌ها پدیدار می‌شوند.

در این پروژه، اعتبارسنجی با اتکا به \textbf{بنچمارک‌های تحلیلی حقیقتِ پایه (\lr{Ground-Truth Analytical Benchmarks})}، آزمون‌های آماری، تست‌های واحد، تست‌های یکپارچه‌سازی، و سناریوهای کامل اجرایی صورت گرفته است. هدف از این ترکیب آن بوده است که از یک سو، اجزای منفرد به‌صورت ایزوله کنترل شوند، و از سوی دیگر، عملکرد کل خط لوله نیز در سناریوهای واقع‌گرایانه سنجیده شود.

\subsection*{۱. یکپارچگی لایه فیزیکی}

نخستین و بنیادی‌ترین لایه‌ای که باید اعتبارسنجی شود، لایهٔ فیزیکی است. در این لایه، منبع نوری، کانال انتقال و آشکارسازها مجموعه‌ای از رخدادها را تولید می‌کنند که تمام مراحل بعدی بر آن‌ها استوار است. اگر این لایه دچار انحراف مفهومی یا عددی باشد، حتی دقیق‌ترین اثبات امنیتی نیز بر ورودی نادرست سوار خواهد شد. به همین دلیل، بخش مهمی از فرایند توسعه به اطمینان از صحت این لایه اختصاص یافت.

در سطح منبع، آمار نشر پالس‌ها برای مدل‌های مورد استفاده—از جمله توزیع پواسونی و در صورت نیاز مدل‌های حرارتی—در برابر روابط تحلیلی شناخته‌شده بررسی شد. در سطح کانال، مدل تضعیف فیبر نوری با تکیه بر قانون بیر–لمبرت و پارامترهای استاندارد انتقال بررسی گردید. در سطح آشکارسازی، رفتار نویز شمارش تاریک، احتمال کلیک، زمان مردهٔ غیرفلج‌شونده (\lr{Non-paralyzable Dead Time}) و آثار وابسته به آن با مدل‌های مرجع تطبیق داده شد.

برای آنکه این تطبیق تنها به مقایسهٔ بصری یا تقریبی محدود نماند، از چارچوب‌های آماری صریح استفاده شد. به‌طور خاص، با به‌کارگیری یک آزمون آماری \(Z\) دوجمله‌ای و اتخاذ سطح معنی‌داری سخت‌گیرانهٔ \(\alpha = 10^{-4}\)، بررسی شد که کمیت‌هایی مانند بهرهٔ کل یا \(Q_\mu\)، و نرخ خطای بیت کوانتومی (\lr{QBER})، در فواصل و سناریوهای مختلف، در تلورانس تعریف‌شده با پیش‌بینی‌های نظری هم‌خوان باشند. نتیجهٔ این فرایند نشان داد که موتور فیزیکی، در محدوده‌های عملی مورد مطالعه، رفتار مورد انتظار را با دقت قابل قبول بازتولید می‌کند. این موضوع اهمیت زیادی دارد، زیرا نشان می‌دهد انحرافات مشاهده‌شده در مراحل بعد، در صورت وقوع، احتمالاً ناشی از بخش‌های دیگر زنجیره‌اند، نه از لایهٔ پایهٔ فیزیکی.

\subsection*{۲. سازگارسازی چند‌اثباتی رمزنگاری}

یکی از چالش‌های مهم در توسعهٔ فریم‌ورک‌های QKD آن است که بسیاری از پیاده‌سازی‌ها از همان ابتدا حول یک اثبات امنیتی خاص طراحی می‌شوند. در چنین وضعیتی، مسیر داده، ساختار آماری، توابع کمکی و حتی نام‌گذاری متغیرها به‌شکلی ناخواسته به منطق آن اثبات خاص وابسته می‌شود. نتیجه آن است که افزودن اثباتی دیگر، به‌جای آنکه صرفاً توسعهٔ یک ماژول جدید باشد، به بازنویسی بخش‌هایی از کل خط لوله نیاز پیدا می‌کند. این وابستگی پنهان، یکی از موانع اصلی برای مقایسهٔ دقیق اثبات‌های امنیتی در محیط‌های فیزیکی مشترک است.

در نسخه‌های اولیه‌تر این پروژه نیز خط لولهٔ اجرایی تا حدی به اثبات کلید محدود \texttt{Lim2014} گره خورده بود. یکی از مهم‌ترین نقاط عطف در بلوغ معماری، معرفی یک \textbf{لایهٔ سازگارساز (\lr{Adapter Layer})} در اسکریپت اجرایی اصلی بود. این لایه نقش مترجم ساختاری را ایفا می‌کند: داده‌ها و آمارهای خام یا نیمه‌پردازش‌شدهٔ حاصل از شبیه‌سازی را دریافت می‌کند و آن‌ها را به امضاها، ساختارها و قراردادهای مورد انتظار هر اثبات امنیتی تبدیل می‌نماید. به این ترتیب، منطق فیزیکی و منطق اثباتی از یکدیگر جدا می‌شوند، در حالی که همچنان از یک زبان دادهٔ مشترک بهره می‌برند.

بر پایهٔ این بازآرایی معماری، فریم‌ورک توانست چند خانوادهٔ اثباتی را با موفقیت پشتیبانی و اجرا کند:
\begin{itemize}
	\item \textbf{2014 Lim}: به‌عنوان یکی از چارچوب‌های مهم \texttt{Finite-key}، این اثبات توانست در مقیاس‌هایی نظیر \(N = 10^7\)، نرخ‌های کلید محدود بهینه و معناداری تولید کند و به‌عنوان مرجع اصلی در بسیاری از مقایسه‌ها به‌کار رود.
	\item \textbf{2005 Wang}: این اثبات نرخ‌های کلیدی متمایزتر و در مواردی محافظه‌کارانه‌تر تولید کرد و نشان داد که تفاوت در روش کران‌بندی آماری و مفروضات تحلیلی چگونه مستقیماً بر خروجی نهایی اثر می‌گذارد.
	\item \textbf{2005 Ma}: اجرای این اثبات، به‌ویژه در نواحی‌ای که مفروضات گاوسی یا کفایت آماری به‌خوبی برقرار نیستند، به‌درستی رفتار محافظه‌کارانه نشان داد و در برخی رژیم‌ها منجر به نرخ کلید صفر شد. همین رفتار، خود نشانه‌ای از صحت سازگارسازی است؛ زیرا سامانه در این حالت تلاش نکرده است به‌صورت مصنوعی نرخ مثبت تولید کند.
	\item \textbf{Proof Tight}: این شاخه با بهره‌گیری از ساختار Lim2014 و در عین حال اعمال تخمین‌های بهینه‌تر برای برخی کمیت‌ها، رفتار متفاوتی را در تحلیل امنیتی ارائه داد و نشان داد که فریم‌ورک آمادهٔ مقایسهٔ نسخه‌های سخت‌گیرانه‌تر یا دقیق‌تر از اثبات‌های مرجع است.
\end{itemize}

اهمیت این دستاورد تنها در «اجرا شدن» چند اثبات خلاصه نمی‌شود. ارزش اصلی آن در این است که پژوهشگر اکنون می‌تواند دقیقاً یک سناریوی فیزیکی واحد را تحت چند لنز امنیتی مختلف مشاهده و تحلیل کند. این قابلیت، برای بررسی اثر انتخاب اثبات بر نرخ کلید، مقایسهٔ محافظه‌کاری یا تهاجمی بودن کران‌ها، و مطالعهٔ حساسیت خروجی به فروض آماری، اهمیت پژوهشی بالایی دارد.

\subsection*{۳. بهینه‌سازی پارامترها با برنامه‌ریزی خطی (LP)}

سامانه‌های QKD ذاتاً به پارامترهای عملیاتی خود حساس‌اند. شدت سیگنال (\(\mu\))، شدت‌های فریب (\(\nu\))، احتمال انتخاب هر شدت، نسبت پایه‌ها، و گاه حتی تنظیمات اجرایی آشکارساز یا پنجره‌های گیتینگ، همگی می‌توانند به‌طور مستقیم و غیرمستقیم بر نرخ کلید امن، QBER، بازده و برد نهایی سامانه اثر بگذارند. در بسیاری از مطالعات ساده، این پارامترها یا به‌صورت دستی تنظیم می‌شوند یا از مقادیر مرجع در مقاله‌های پیشین اقتباس می‌گردند. با این حال، چنین رویکردی معمولاً برای دستیابی به بهترین عملکرد در یک سناریوی مشخص کافی نیست.

در این فریم‌ورک، یک لایهٔ بهینه‌سازی برای جست‌وجوی نظام‌مند در فضای پارامترها توسعه یافته است. این لایه با استفاده از حلگرهای عددی مانند \lr{Nelder-Mead} و همچنین ابزارهای مبتنی بر \textbf{برنامه‌ریزی خطی (\lr{LP})}، می‌تواند فضای شدت‌ها و احتمال‌ها را پیمایش کرده و پیکربندی‌هایی را بیابد که نرخ کلید امن را در شرایط معین بیشینه می‌کنند. اعتبارسنجی این بخش صرفاً به بررسی همگرایی عددی محدود نشد، بلکه خروجی‌های بهینه‌سازی با سناریوهای مرجع و مقادیر غیر‌بهینه مقایسه شدند تا اثر واقعی بهینه‌سازی بر عملکرد نهایی سنجیده شود.

نتایج این بررسی‌ها نشان دادند که حلگر LP و لایهٔ بهینه‌سازی نه‌تنها به‌درستی در فضای پارامتر حرکت می‌کنند، بلکه در بسیاری از فواصل عملی—به‌ویژه در بازهٔ ۰ تا ۱۰۰ کیلومتر—توانسته‌اند پیکربندی‌هایی بیابند که به \textbf{افزایش بسیار معنادار در نرخ کلید امن} نسبت به مقادیر ثابت و سخت‌کدگذاری‌شده منجر می‌شوند. این نتیجه از منظر مهندسی اهمیت زیادی دارد، زیرا نشان می‌دهد بخشی از شکاف عملکردی میان یک سامانهٔ «قابل اجرا» و یک سامانهٔ «خوب تنظیم‌شده» را می‌توان به‌صورت محاسباتی و سیستماتیک کاهش داد. از منظر پژوهشی نیز، وجود این لایه باعث می‌شود که مقایسهٔ اثبات‌ها یا مدل‌ها بر اساس پارامترهای منصفانه‌تر و نزدیک‌تر به بهینه انجام شود، نه صرفاً بر اساس انتخاب‌های دستی.

\subsection*{۴. توپولوژی شبکه و رله‌های قابل اعتماد}

بخش بزرگی از ادبیات کلاسیک QKD بر لینک‌های نقطه‌به‌نقطه متمرکز است؛ حال آنکه در کاربردهای عملی، به‌ویژه زمانی که هدف فراتر رفتن از اتصال مستقیم دو نقطه باشد، مسئلهٔ شبکه‌سازی کوانتومی مطرح می‌شود. در چنین محیطی، دیگر تنها نرخ کلید یک لینک خاص مهم نیست، بلکه نحوهٔ توزیع و مصرف کلیدها در سراسر توپولوژی، مدیریت لینک‌های میانی، تاب‌آوری در برابر خرابی، و منطق مسیریابی نیز به مسئله افزوده می‌شود. به همین دلیل، یکی از ابعاد مهم اعتبارسنجی این پروژه، توسعه و بررسی یک لایهٔ شبکه بر بستر موتور فیزیکی و امنیتی بود.

اوج این مرحله از اعتبارسنجی در شبیه‌سازی یک \textbf{توپولوژی ۷ گرهی} حاصل شد؛ توپولوژی‌ای که شامل فرستنده‌ها، گیرنده‌ها و رله‌های قابل اعتماد (\lr{Trusted Relays}) بوده و امکان بررسی سناریوهای چندمرحله‌ای را فراهم می‌کرد. در این سناریو، موتور شبکه باید علاوه بر مدیریت تولید کلید روی لینک‌های منفرد، بتواند از این کلیدها در سطح بین‌گرهی نیز استفاده کند. این مسئله شامل نگهداری و به‌روزرسانی \textit{استخرهای کلید} (\lr{Key Pools})، ارزیابی هزینهٔ مسیرها، تصمیم‌گیری برای انتخاب مسیر مناسب، و در صورت لزوم، رله‌گذاری مرحله‌به‌مرحلهٔ کلید میان گره‌هاست.

اعتبارسنجی این لایه نشان داد که سامانه توانایی انجام موارد زیر را دارد:
\begin{itemize}
	\item اجرای دقیق منطق \textbf{کوتاه‌ترین مسیر اول (\lr{SPF+})} بر اساس هزینهٔ لینک و وضعیت استخر کلیدها؛
	\item انجام \textbf{رله‌گذاری مرحله‌به‌مرحله (\lr{Hop-by-Hop Relay})} و مدیریت ترکیب کلید در مسیرهای چندگامی؛
	\item نمایش \textbf{Failover پویا}، به این معنا که در صورت حذف یا قطع عمدی یک لینک اصلی، سامانه بتواند مسیر جایگزین مناسبی را شناسایی و سرویس توزیع کلید را از مسیر جدید ادامه دهد؛
	\item تغذیه و بازیابی خودکار استخرهای کلید برای جلوگیری از فروپاشی عملکرد شبکه در بارهای متغیر.
\end{itemize}

این سطح از اعتبارسنجی نشان می‌دهد که فریم‌ورک حاضر صرفاً ابزاری برای تولید چند شاخص فیزیکی نیست، بلکه می‌تواند در مطالعهٔ معماری‌های گسترده‌تر QKD و رفتار شبکه‌ای نیز نقش ایفا کند. این ویژگی به‌ویژه برای پژوهش‌هایی که به استقرار شبکه‌های کوانتومی کاربردی، مدیریت کلید در مقیاس چندگرهی، یا ارزیابی تاب‌آوری مسیرها علاقه‌مندند، اهمیت ویژه دارد.

\section*{تاب‌آوری مهندسی}

هر پروژهٔ نرم‌افزاری که به‌طور جدی در معرض توسعه، بازطراحی، افزودن قابلیت و اعتبارسنجی مستمر قرار گیرد، دیر یا زود با موارد لبه، تضادهای معماری، یا ناهماهنگی‌های پنهانی روبه‌رو می‌شود. در سامانه‌ای مانند این فریم‌ورک، که در آن لایه‌های فیزیک، آمار، امنیت، همزمانی و شبکه در تعامل‌اند، چنین مواردی نه استثنا، بلکه بخشی طبیعی از مسیر بلوغ مهندسی هستند. تفاوت اصلی میان یک نمونهٔ آزمایشگاهی و یک ابزار پژوهشی قابل اتکا، تا حد زیادی در نحوهٔ مواجهه با همین جزئیات نهفته است.

در طول توسعه و اعتبارسنجی این فریم‌ورک، چندین مورد لبه شناسایی و اصلاح شدند که رفع آن‌ها نقشی مهم در افزایش تاب‌آوری سیستم ایفا کرد. بخشی از این اصلاحات به \textbf{تضادهای ارث‌بری} در امضاهای متدهای برخی اثبات‌ها، از جمله شاخهٔ \texttt{Tight}، مربوط می‌شد. اگرچه این دسته از مسائل ممکن است در نگاه اول صرفاً به سازگاری کدنویسی مربوط به نظر برسند، اما در عمل می‌توانند منجر به اجرای نادرست متدها، تفسیر غلط داده‌های ورودی، یا ناسازگاری خاموش در مسیرهای اجرایی شوند. رفع این تعارض‌ها باعث شد قراردادهای میان کلاس‌های پایه و مشتق‌شده شفاف‌تر و پایدارتر شوند.

بخش دیگری از اصلاحات به \textbf{سریال‌سازی و بازیابی وضعیت}، به‌ویژه در سطح آشکارساز، مربوط بود. هنگامی که یک سامانهٔ شبیه‌سازی قرار است حالت اجزای خود را ذخیره و بازیابی کند، عدم کنترل دقیق نسخهٔ طرح‌واره (\lr{Schema Version}) می‌تواند به برش خاموش داده‌ها، جابه‌جایی فیلدها، یا بازیابی ناقص وضعیت منجر شود. به همین دلیل، بررسی‌های سخت‌گیرانه‌تری برای سازگاری نسخهٔ طرح‌واره و درستی بازخوانی داده‌ها اعمال شد. این موضوع برای سناریوهایی که در آن‌ها تداوم وضعیت آشکارساز، carry-over رویدادها یا ادامهٔ اجرای دسته‌ای اهمیت دارد، حیاتی است.

همچنین، برخی اصلاحات به \textbf{چینش فیلدهای تغییرناپذیر در ساختارهای داده} و نحوهٔ پذیرش پارامترهای امنیتی چندگانه بازمی‌گشت. در ساختارهای مبتنی بر \textit{dataclass}، کوچک‌ترین ناسازگاری در ترتیب فیلدها یا قرارداد سازنده می‌تواند در نسخه‌های مختلف یا هنگام گسترش مدل، به رفتارهای شکننده منتهی شود. بازتنظیم این بخش‌ها، انعطاف‌پذیری سیستم در افزودن پارامترهای جدید را افزایش داد و احتمال بروز خطاهای پنهان در ساخت اشیا را کاهش داد.

اهمیت این اصلاحات در آن است که نشان می‌دهند تاب‌آوری یک فریم‌ورک، تنها از طریق داشتن مدل‌های نظری خوب حاصل نمی‌شود. یک ابزار پژوهشی باید در برابر ورودی‌های مرزی، سناریوهای غیرمعمول، توسعهٔ تدریجی و تنوع مسیرهای اجرا نیز پایدار باشد. از این منظر، تاب‌آوری مهندسی بخشی از اعتبار علمی پروژه است؛ زیرا نتایجی که بر بستر نرم‌افزار شکننده تولید شوند، حتی اگر از نظر فرمول‌بندی صحیح باشند، در سطح اعتمادپذیری پژوهشی با محدودیت مواجه خواهند بود.

\section*{راهنمای مطالعه این مستند}

این مستند با این فرض تدوین شده است که خواننده ممکن است از یکی از چند مسیر متفاوت وارد آن شود: از مسیر نظریهٔ \texttt{QKD}، از مسیر پیاده‌سازی نرم‌افزار، از مسیر اعتبارسنجی پژوهشی، یا از مسیر کاربردهای شبکه‌ای. به همین دلیل، ساختار سند به‌گونه‌ای طراحی شده که هم بتوان آن را به‌صورت ترتیبی از ابتدا تا انتها مطالعه کرد و هم بتوان از آن به‌عنوان یک مرجع موضوع‌محور استفاده نمود.

به‌طور کلی، این سند در چند بخش اصلی سازمان‌دهی شده است:

\begin{itemize}
	\item \textbf{بخش اول: معماری فریم‌ورک و فیزیک پایه} \\
	در این بخش، مبانی مفهومی و محاسباتی لایهٔ فیزیکی معرفی می‌شوند. مدل‌های منبع نوری، کانال، تلفات، نویز، و آشکارسازها در این قسمت تشریح می‌گردند و روشن می‌شود که هر یک چگونه در قالب ساختار نرم‌افزاری پروژه پیاده‌سازی شده‌اند.
	
	\item \textbf{بخش دوم: اثبات‌های امنیتی رمزنگاری} \\
	این بخش به منطق امنیتی فریم‌ورک اختصاص دارد. در آن، چارچوب‌های پشتیبانی‌شده برای تخمین پارامتر، کران‌بندی آماری و محاسبهٔ نرخ کلید امن ارائه می‌شوند. افزون بر فرمول‌بندی نظری، ارتباط این اثبات‌ها با رابط‌های نرم‌افزاری و داده‌های ورودی نیز توضیح داده می‌شود.
	
	\item \textbf{بخش سوم: خط لوله اجرایی و لایه شبکه} \\
	در این بخش، مسیر کامل پردازش داده از مرحلهٔ اجرای شبیه‌سازی تا تولید خروجی و همچنین منطق موتور توزیع، کارگرها، هماهنگی موازی و قابلیت‌های شبکه‌ای تشریح می‌شود. این بخش برای فهم رفتار سامانه در مقیاس بزرگ و سناریوهای چندگرهی اهمیت ویژه دارد.
	
	\item \textbf{بخش چهارم: اعتبارسنجی و تست} \\
	در این قسمت، بنچمارک‌های تحلیلی، روش‌های آزمون، نتایج اعتبارسنجی، و سناریوهای تستی که برای راستی‌آزمایی فریم‌ورک به‌کار رفته‌اند، ارائه می‌شود. این بخش برای خوانندگانی که به اعتمادپذیری نتایج و بازتولیدپذیری علاقه‌مندند، اهمیت محوری دارد.
\end{itemize}

در مجموع، هدف از این ساختاربندی آن است که مستند حاضر صرفاً توصیفی از یک کدبیس نباشد، بلکه پلی میان \textbf{مبانی نظری QKD}، \textbf{مهندسی دقیق نرم‌افزار} و \textbf{اعتبارسنجی پژوهشی} ایجاد کند. به بیان دیگر، این سند باید هم برای کسی که می‌خواهد منطق فیزیکی و امنیتی فریم‌ورک را درک کند مفید باشد، هم برای توسعه‌دهنده‌ای که می‌خواهد آن را گسترش دهد، و هم برای پژوهشگری که می‌خواهد به نتایج آن استناد کند یا آن‌ها را بازتولید نماید.

امید آن است که این فریم‌ورک بتواند به‌عنوان ابزاری \textbf{قابل‌اعتماد، شفاف، توسعه‌پذیر و بازتولیدپذیر} در اختیار پژوهشگران حوزهٔ رمزنگاری کوانتومی قرار گیرد؛ ابزاری که نه‌تنها برای تحلیل یک سناریوی منفرد، بلکه برای ساختن و ارزیابی نسل بعدی مطالعات در حوزهٔ استقرار سامانه‌ها و شبکه‌های کوانتومی کاربردی مفید باشد. اگر این پروژه بتواند حتی بخشی از مسیر میان صورت‌بندی‌های نظری QKD و سامانه‌های عملی و مهندسی‌شدهٔ آن را روشن‌تر کند، به هدف اصلی خود نزدیک شده است.
\end{sloppypar}
\vspace{2cm}
\begin{flushleft}
	\textit{رضا آرام جو}\\
	\textit{\today}
\end{flushleft}
